\chapter[\emj{🌲}~Fenwick Tree / BIT (Binary Indexed Tree)]{\emj{🌲}~Fenwick Tree / BIT (Binary Indexed Tree)}

\begin{dsaquees}

Una versión súper compacta del Segment Tree para sumas de prefijo 💰.
Imagina que quieres el saldo acumulado hasta el día 7, pero los montos
cambian a diario.

El Fenwick Tree guarda sumas parciales inteligentes: cada posición del
array "cubre" un bloque de tamaño potencia de 2, definido por su bit
menos significativo. Para sumar hasta i, saltas hacia atrás quitando ese
bit; para actualizar, saltas hacia adelante sumándolo. Ambos recorridos
son $O(\log n)$ y el código cabe en cuatro líneas.

\end{dsaquees}

\begin{dsaotraforma}

Un Fenwick Tree (o Binary Indexed Tree, BIT) es una estructura que
calcula sumas de prefijo y actualizaciones de punto en $O(\log n)$,
usando solo un array de tamaño $n+1$ (mucho menos que el $4n$ del Segment
Tree).

Su magia es el operador \verb|i & (-i)|, que aísla el bit menos
significativo de $i$ y define el tamaño del bloque que cada índice
resume:
\begin{itemize}
\item \textbf{Query prefijo(i)}: parte de $i$ y resta el bit menos
      significativo hasta llegar a 0, acumulando.
\item \textbf{Update(i, delta)}: parte de $i$ y SUMA el bit menos
      significativo hasta pasar $n$, propagando el cambio.
\end{itemize}

Una suma de rango $[l, r]$ se obtiene como $\mathrm{prefijo}(r) -
\mathrm{prefijo}(l-1)$.

\end{dsaotraforma}

\begin{dsadiagrama}
\begin{verbatim}
Índices (1-based) y qué bloque cubre cada uno (por lowbit):

  i:     1    2    3    4    5    6    7    8
  cubre: [1] [1-2][3] [1-4][5] [5-6][7] [1-8]

lowbit(i) = i & (-i)   (bit menos significativo)
  lowbit(6)=2 -> el 6 cubre 2 elementos: [5,6]
  lowbit(8)=8 -> el 8 cubre 8 elementos: [1..8]

query prefijo(7): 7 -> 6 -> 4 -> 0   (resta lowbit)
                  suma tree[7]+tree[6]+tree[4]
update(5,+3):     5 -> 6 -> 8 -> fin  (suma lowbit)
\end{verbatim}
\end{dsadiagrama}

\begin{dsacomofunciona}
\begin{verbatim}
QUÉ RANGO cubre cada índice (largo = lowbit)  y los SALTOS

  i:      1    2    3    4    5    6    7    8
          │   ┌┴┐   │  ┌─┴──┐  │  ┌┴┐   │ ┌──┴────┐
  cubre: [1] [1-2] [3][1 - 4][5][5-6] [7][1  -  8]

QUERY prefijo(7) = suma de 1..7   (resta lowbit: baja)
   7 ──lowbit1─► 6 ──lowbit2─► 4 ──lowbit4─► 0 (fin)
   suma = tree[7] + tree[6] + tree[4]
        = [7]   + [5-6]  + [1-4]   = cubre 1..7  OK

UPDATE(3, +5)  (suma lowbit: sube)
   3 ──lowbit1─► 4 ──lowbit4─► 8 ──► fin (>n)
   tree[3] += 5 ; tree[4] += 5 ; tree[8] += 5
   (todos los bloques que "contienen" al índice 3 se enteran)

Cada salto quita/suma un bit → a lo más log(n) pasos.
\end{verbatim}
\end{dsacomofunciona}

\begin{lstlisting}
CÓMO FUNCIONA (1-indexado)
- lowbit(i) = i & (-i) aísla el bit menos significativo.
- update(i, d): mientras i <= n: tree[i] += d; i += lowbit(i).
- query(i): suma = 0; mientras i > 0: suma += tree[i]; i -= lowbit(i).
- rango(l, r) = query(r) - query(l-1).

📝 PSEUDOCÓDIGO
──────────────────────────────────────────────────────────────

función update(i, delta):
    MIENTRAS i <= n:
        tree[i] += delta
        i += i & (-i)         // avanza al siguiente bloque

función query(i):             // suma de 1..i
    suma = 0
    MIENTRAS i > 0:
        suma += tree[i]
        i -= i & (-i)         // retrocede quitando el bit
    return suma

💻 CÓDIGO C++
──────────────────────────────────────────────────────────────

#include <vector>
using namespace std;

struct Fenwick {
    int n; vector<long long> tree;
    Fenwick(int n) : n(n), tree(n + 1, 0) {}   // 1-indexado

    void update(int i, long long delta) {
        for (; i <= n; i += i & (-i)) tree[i] += delta;
    }
    long long query(int i) {                   // suma de 1..i
        long long s = 0;
        for (; i > 0; i -= i & (-i)) s += tree[i];
        return s;
    }
    long long rango(int l, int r) {            // suma de l..r
        return query(r) - query(l - 1);
    }
};

⏱️ COMPLEJIDAD (Big-O)
──────────────────────────────────────────────────────────────

  Operación       │ Tiempo     │ Espacio
  ────────────────┼────────────┼──────────
  Update de punto │ O(log n)   │ O(n)
  Query prefijo   │ O(log n)   │ O(n)
  Construir       │ O(n log n) │ O(n)

* Solo array de tamaño n+1: mucho más ligero que el Segment Tree.
\end{lstlisting}

\begin{dsaentrevistas}

✅ Fenwick vs Segment Tree: si solo necesitas sumas de prefijo con
   updates de punto, Fenwick gana (código corto, menos memoria, más
   rápido en la práctica). Para mínimo/máximo o updates de rango, ve al
   Segment Tree.

💡 "Count of Smaller Numbers After Self" (LeetCode 315): se resuelve con
   un BIT sobre valores comprimidos. Clásico truco de conteo de
   inversiones.

⚠️ Es 1-indexado: el índice 0 no se usa. Si tu array es 0-based, suma 1
   al indexar o te quedas en un bucle infinito con \verb|i & (-i)|.

🔑 Conteo de inversiones: recorre de derecha a izquierda y consulta
   cuántos elementos menores ya insertaste. Es EL uso estrella del BIT en
   entrevistas.

\end{dsaentrevistas}

\begin{dsaresumen}

• BIT: sumas de prefijo + update de punto en O(log n).
• Usa \verb|i & (-i)| (bit menos significativo) para saltar bloques.
• query: resta el bit; update: suma el bit.
• Rango [l,r] = prefijo(r) − prefijo(l−1).
• Es 1-indexado; array de tamaño n+1 (más ligero que Segment Tree).
• Estrella para contar inversiones y "smaller after self".
════════════════════════════════════════════════════════════════

\end{dsaresumen}
